\section{From a return declaration to the next annual opening}
\label{sec:company-return-reference}

P's recovered stock gives the language a concrete implementation problem.
The return contract, the legal book and the tax calculation each remember
something different about the same two units. The contract records that the
customer's rights have ended. The stock record keeps the recovered carrying
amount until sale. The tax calculation remembers an adjustment arising from
expense recognized before the year began. A company program must preserve
all three without posting the opening balances again.

The finite reference now connects these records through the complete retained
2025 P trace. It implements return settlement, whole-lot sale, P's selected
current and deferred tax method, legal statements and annual carry. Other
corporate events continue to use the checked retained methods. This division
makes the implementation comparison precise: the new return and tax
calculations can be compared with the earlier company case, while the
surrounding methods supply the same corporate transactions as before.

\paragraph{An author declares facts before choosing an event.}
The versioned \code{ReturnModule}\allowbreak\code{AuthoringDraft/1} describes a component of
one legal book. It identifies the source packet and opening state, entity,
currency, period and account roles. Its contract memory includes sold and
expected-return units, unit refund and cost, estimated recovery costs, and
outstanding refund and recovery balances. The recovered-stock declarations
keep quantities, carrying amounts and tax bases separately. Tax declarations
state the selected timing method, rate and required findings.

These declarations allocate existing accounts. For example, a refund balance
of 40 must agree with the liability in P's opening book; it is not an
instruction to create another liability. The deferred asset and liability
must likewise agree with their declared opening allocations. The remaining
company balances come from the shared opening, rather than from a second
component wallet.

Typed records and text declarations contain the same fields. The text form
uses named declarations for the component, lots, tax premises and dated
settlement or sale. It contains no executable Python, external include or
opaque JSON body. Both forms preserve canonical exact rational amounts, strict
Boolean findings and source references. Rendering a typed declaration as text
and reading it again must reproduce the same value. This is a matched
authoring comparison with the typed-library baseline, not evidence that text
is preferable for every user.

Unknown values are explicit. Suppose the tax rate has not been established
but the author also supplies a refund liability of 41 while two expected
returns at 20 require 40. The contradiction is already decidable and the
record is invalid. The unknown tax rate cannot conceal it. If the amounts
agree but a required recovery finding is false, the selected method is
unsupported. If the required finding is unknown, the affected calculation
remains unresolved. Arithmetic checks inspect their own dependencies rather
than replacing unknowns with zero or postponing every known check.

Source correspondence is another question. A syntactically valid source
pointer may name the wrong disclosure. The retained example therefore checks
the entire declared projection against P's original packet, including owner,
dates, amounts and opening identity. Changing P to Northstar fails that
correspondence even though both names are valid text. A correctly bound local
packet still needs separate evidence for external authenticity and any
historical admission claim.

\paragraph{One event changes one shared book.}
The reference state combines P's account balances with return memory,
recovered lots, dated tax adjustments, classified gross cash and lifetime
event identities. Every account belongs to one financial-statement class;
every temporary account enters profit. A proposed return first checks its
component, quantities, findings and date. It then derives the settlement
journal already shown in Section~\ref{sec:company-module-language}, creates
the separate recovered lot and closes the return rights. The surrounding
inventory allocation stays intact.

The proposal identifies the complete preceding state, the exact event and
the selected method bundle. Applying it repeats the calculation with the
trusted local method and compares the journal, gross cash and complete
after-state. A changed cash amount cannot become valid merely by supplying a
new hash. A stale proposal, repeated settlement or sale of an already disposed
lot is refused. This is an in-process reference for checking a transition;
durable publication and concurrent writers require the separate storage design.

Cash sufficiency is checked after deriving the complete proposed journal.
If P cannot pay the refund and handling costs, the event returns a funding
failure and leaves the preceding book intact. It does not create financing.
Missing tax evidence has a different consequence: the valid non-tax events
can remain as an exact dated prefix, but the tax-dependent annual statement
and next-year opening remain unavailable. Thus a useful partial history
survives without being presented as a complete year.

The retained adapter has an explicit, narrower role. It recomputes the other
corporate events from their original state and checked methods, verifies that
their P opening balances agree with the shared book, and supplies their
journals and cash origins. Saved expected journals are comparison evidence,
not instructions to the executor. S's book and the group ownership state
remain in that retained method. Replaying them checks compatibility; it does
not claim a second independent implementation of those same corporate rules.

\paragraph{Tax accumulates within a year and resets at its boundary.}
The settlement records its taxable-minus-book adjustment of $-24$ with its
event identity and date. At either tax close, let $B_y(t)$ be P's cumulative
pretax book income from the start of year $y$ through cutoff $t$, let
$A_y(t)$ be the sum of the selected method's dated tax adjustments in that
interval, and let $C_y(t^-)$ be current tax already recognized in the year.
Under the retained constant rate $r=1/4$, the additional current tax is
\begin{equation}
 \Delta C_y(t)=r\{B_y(t)+A_y(t)\}-C_y(t^-).
 \label{eq:return-reference-current-tax}
\end{equation}
The final subtraction prevents an interim tax charge from being paid again.
At the first close, $B_y=56$, $A_y=-24$ and $C_y(t^-)=0$, giving 8. At the
annual close, $B_y=66$ and the same cumulative adjustment remains $-24$;
subtracting the existing 8 gives the remaining 2.5. Negative additional tax
would require a refund or other tax method that this finite reference does
not supply. The opening deferred balances reverse separately, producing
deferred expense 6.

The annual legal statements reconcile profit 49.5, cash 491.5 and equity
671.5. Their operating cash of 101.5 agrees with both the gross payments and
the indirect bridge in equation~\eqref{eq:module-return-cash}. In particular,
the opening interest payment remains a cash adjustment even when current
finance expense is zero. Legal receipts from S stay in P's cash flow; they
are eliminated only by the separate group method.

At year end, the reference transfers profit to retained earnings and clears
the income and period cash accumulators. It retains the source context,
settled return component, disposed-lot identity and lifetime event history.
The old tax adjustment remains available to explain 2025, but its date is
outside $A_{2026}(t)$. In a deliberately inactive following year, equation~
\eqref{eq:return-reference-current-tax} therefore gives zero additional tax
and the deferred reversal is not repeated. This checks annual memory; it is
not a prediction of an inactive business.

The complete reference compares every P account and gross cash origin after
all thirteen dated transitions, then checks the annual statements and carry.
It also preserves the existing six company-authoring paths. The remaining
language task is a composed-company declaration that makes each selected
method, state allocation and unresolved dependency explicit. General changes
in return estimates, partial lot sales, other tax treatments and production
execution still need their own design and conformance cases.

The matched authoring and execution records are under
\path{data/accounting-master/}:
\begin{quote}\small
\path{company-return-schema-20261023}\\
\path{company-return-execution-20261024}
\end{quote}
They retain
the original source comparisons and failure cases. The resulting methods are
conditional finite references; independent implementation, domain review and
human acceptance of the governing design remain pending.
